\section{7.1 n-step TD: MC와 TD(0) 사이의 다이얼}
\begin{frame}{복습: 두 극단 --- MC와 TD(0)}
  \begin{itemize}
    \item MC의 목표값: 에피소드 \textbf{끝까지의} 실제 보상 합(실제 리턴)
      $G_t$ --- 추정치 없음.
    \item TD(0)의 목표값: $r + \gamma V(s')$ --- 실제 보상 $1$개 + 나머지
      추정치.
    \item 나란히 쓰면 \textbf{같은 구조의 다른 길이}: 차이는 ``실제 보상을
      몇 개 쓸 것인가''가 $1$개와 ``끝까지 전부''인 것.
  \end{itemize}
  \vskip0.8em
  \kq{1과 ``끝까지'' 사이의 아무 정수 $n$도 가능해야 하지 않나 ---
  $n$개의 실제 보상 + 나머지는 추정치, 라는 목표값을 쓰자. 이것이
  \kb{n-step TD}.}
\end{frame}
